% Lemma: k=3 base case for Theorem 7.5 (Non-square discriminant)
% Generated and verified by verify_k3base.py

\begin{lemma}\label{lem:k3base}
For $k=3$, the $24$ all-one triples $($up to the ordering
convention $\eps_a < \eps_b)$ yield exactly two equivalence classes
of Gaussian elements $z_+ \in \Q(i)[t_1, t_3]$, under the equivalence
$z_+ \sim u \cdot z_+({\pm t_1, t_3})$ for $u \in \{{\pm 1, \pm i}\}$.
A representative of each class is
\begin{align}
  z_+ &= -(t_3 - i)^2\,(t_1^2 - 1)
         \;+\; 2\bigl(i(t_3^2 - 1) - 6t_3\bigr)\,t_1,
         \label{eq:zplus-class0} \\
  z_+' &= \phantom{-}(t_3 + i)^2\,(t_1^2 - 1)
         \;-\; 2\bigl(i(t_3^2 - 1) + 6t_3\bigr)\,t_1,
         \label{eq:zplus-class1}
\end{align}
and $z_+'(t_1, t_3) = -\overline{z_+(-t_1, t_3)}$,
where the bar denotes conjugation in~$\Q(i)/\Q$.
Both $z_+$ and $z_+'$ are irreducible over~$\Q(i)$, and in particular
neither is a Gaussian square.
\end{lemma}

\begin{proof}
The enumeration of the $24$ triples and the reduction to two equivalence
classes is a finite computation (verified in the accompanying script
\texttt{verify\_k3base.py}).

It suffices to prove that $z_+$ in~\eqref{eq:zplus-class0} is irreducible
over~$\Q(i)$; the irreducibility of~$z_+'$ then follows since
$z_+' = -\overline{z_+(-t_1, t_3)}$ and the map
$f(t_1, t_3) \mapsto -\overline{f(-t_1, t_3)}$ is a ring automorphism
of $\Q(i)[t_1, t_3]$.

As a polynomial in~$t_1$, $z_+$ has degree~$2$ with coefficients
in~$\Q(i)[t_3]$:
\[
  z_+ = a\,t_1^2 + b\,t_1 + c, \qquad
  a = -(t_3 - i)^2, \quad
  b = 2\bigl(i(t_3^2-1) - 6t_3\bigr), \quad
  c = (t_3 - i)^2 = -a.
\]
The discriminant is
\[
  \Delta(t_3) \;=\; b^2 - 4ac \;=\; b^2 + 4a^2 \;=\; -64\,i\,t_3\,(t_3 + i)^2.
\]
Since $\deg_{t_3}(\Delta) = 3$ is odd, $\Delta(t_3)$ cannot be a perfect
square in~$\Q(i)[t_3]$ (a square polynomial has even degree).
Therefore $z_+$ is irreducible over~$\Q(i)(t_3)$, and a~fortiori
over~$\Q(i)$ (by Gauss's lemma, an element of $\Q(i)[t_1, t_3]$ that is
irreducible in $\Q(i)(t_3)[t_1]$ is irreducible in $\Q(i)[t_1, t_3]$,
since $\Q(i)[t_3]$ is a UFD).

In particular, $z_+$ is not a Gaussian square, so
$D_1/16 = z_+\,\bar{z}_+$ is not a perfect square.
\end{proof}

\begin{remark}
The resultant $\Res_{t_1}(z_+,\, \partial z_+/\partial t_1)
= -64\,i\,t_3\,(t_3^2+1)^2$
is a nonzero polynomial in~$t_3$, confirming that $z_+$ and
$\partial z_+/\partial t_1$ are coprime in~$\Q(i)(t_3)[t_1]$
and hence that $z_+$ is squarefree.
The odd-degree discriminant argument is however more direct.
\end{remark}
